\documentclass[11pt]{article}
\usepackage{preamble}
\renewcommand{\answer}[1]{}

\begin{document}

\ladderheading{AP Statistics \textperiodcentered\ Unit 1 \textperiodcentered\ Topic 1.13 \textperiodcentered\ B: 2026-10-02 \textperiodcentered\ E: 2026-10-14}{Which Design Supports Cause?}

\focusbanner{TODAY'S PROBLEM}

Suppose a researcher recruits 30 student volunteers to study whether an energy drink changes reaction time.\par\smallskip \textbf{Design 1:} The first 15 students to sign up receive the energy drink; the last 15 receive water. A tester who knows each drink measures reaction time in milliseconds.\par \textbf{Design 2:} A coin flip assigns each volunteer to the energy drink or water. The water is flavored to match, both drinks are served in identical opaque cups, and neither the volunteer nor the tester knows which drink was assigned.\par
\begin{center}
\textbf{Two energy-drink designs}\par\smallskip
\begin{tabular}{|l|c|c|}
\hline
\textbf{\#} & \textbf{Drink allocation} & \textbf{Tester told drink?} \\ \hline
\textbf{1} & First/last signup & Yes \\ \hline
\textbf{2} & Coin flip & No \\ \hline
\end{tabular}
\end{center}
\begin{firsttakebox}{FIRST TAKE --- before the video (5 min)}
Which design would make the comparison between drinks fairer? Give one reason.
\workspace{0.9in}
\end{firsttakebox}

\begin{callout}{\IconBook\ PRINCIPLES OF EXPERIMENTAL DESIGN (keep on the page)}{calloutgreen}
A well-designed experiment uses \textbf{comparison}, \textbf{random assignment}, \textbf{replication}, and
control of potential confounding variables. A \textbf{control group} supplies a baseline. Random assignment
tends to balance uncontrolled variables so a difference can be attributed to treatment. \textbf{Blinding}
reduces changes caused by expectations or by how responses are measured. Selection still limits generalization.
\end{callout}

\newpage
\textbf{Finish after the video:}\par\smallskip

\dokbadge{1}~\textbf{(a)} Name the experimental units and the response variable, with units.\par
\workspace{0.8in}

\dokbadge{2}~\textbf{(b)} Compare how treatments are assigned in the two designs. Explain how one other feature of Design 2 helps the comparison.\par
\workspace{1.4in}

\dokbadge{3}~\textbf{(c)} Which design gives stronger evidence about whether the drink causes a change? Explain one other reason the groups in Design 1 might react differently. State why the result would not automatically describe all teenagers.\par
\workspace{2.0in}

\begin{sentenceframebox}\raggedright\textbf{Frame:} Design \blank[0.4in] supports cause because \blank[2.7in].\end{sentenceframebox}

\begin{sentenceframebox}\raggedright\textbf{Frame:} The result cannot automatically generalize to \blank[1.5in] because \blank[2in].\end{sentenceframebox}

\begin{summaryexitbox}{EXIT --- turn this in}
One thing the video changed about my first take: \hrulefill\par
\workspace{0.4in}
\end{summaryexitbox}

\end{document}
